\documentclass[11pt]{article}

\usepackage[T1]{fontenc}
\usepackage{amsmath,amssymb,amsthm}
\usepackage[margin=1in]{geometry}
\usepackage[hidelinks,pdfusetitle]{hyperref}

\newtheorem{theorem}{Theorem}
\newtheorem{lemma}[theorem]{Lemma}
\newtheorem{proposition}[theorem]{Proposition}
\theoremstyle{remark}
\newtheorem{remark}[theorem]{Remark}

\newcommand{\Z}{\mathbb Z}
\newcommand{\one}{\mathbf 1}
\newcommand{\E}{\mathbb E}
\newcommand{\Prob}{\mathbb P}
\DeclareMathOperator{\dist}{dist}

\title{A Covariance Separator for an Explicit\\
Dean--Goldsmith Lemma~7 Sample Family}
\author{Xiaojian Zeng\\
PBC School of Finance, Tsinghua University, 100190\\
\texttt{xiaojianzeng@foxmail.com}}
\date{}

\begin{document}
\maketitle

\begin{abstract}
We analyze one explicit family of samples produced by the
bounded-distance-decoding construction in Lemma~7 of Dean and Goldsmith.
For every sufficiently large dimension \(n\), take
\(L=\Z^n\), \(A=I_n\), \(r=k=n^2\), \(\alpha=1/2\),
\(m=3n^5\), any admissible constellation size \(M\), and the target
\(y=2M\one+e_1/4\).  We verify all relevant parameter and
bounded-distance promises.  We then prove that \(N=n^8\) independent
constructed samples are separated from every fixed legal discrete-MIMO
law, uniformly over its secret and every noise scale at most \(1/2\), by
thresholding one empirical covariance ratio.  An explicit negligible
upper bound is given for both classification errors.  The proof controls
the unbounded discrete-Gaussian public vectors, the empirical denominator,
and response noise on the same finite-sample event.  This resolves the
stated explicit family; it is neither an attack on general MIMO decoding
nor a hardness result.
\end{abstract}

\section{Context and statement}

Dean and Goldsmith proposed a reduction from lattice problems to decoding
in a massive-MIMO model \cite{DeanGoldsmith2017}.  Their construction
adapts the lattice and learning-with-errors framework of Regev
\cite{Regev2009}; discrete Gaussian measures and their smoothing
parameter are central to the worst-case/average-case theory of
Micciancio and Regev \cite{MicciancioRegev2007}.  Attacks on related
parameter choices and constructions were subsequently given by Steinfeld
and Sakzad \cite{SteinfeldSakzad2015,SakzadSteinfeld2020}.

Li's proof-level analysis \cite{Li2026} identifies a separate issue in the
Lemma~7 sample constructor.  Proposition~5.1 there proves exact
non-membership in the legal discrete-MIMO family when the BDD coefficient
is outside the constellation or the BDD offset is nonzero.  Exact
non-membership alone does not imply efficient statistical
distinguishability; Remark~5.3 of \cite{Li2026} explicitly leaves this
question dependent on the instance.  The result below closes that
finite-sample gap for one completely specified family.  We make no
priority claim beyond this bounded comparison with the cited sources.

The scope of an explicit certificate matters: resolving one quantified
family need not resolve its ambient problem.  The same discipline appears,
in a different setting, when an exact matrix counterexample is presented
together with the parameter regimes that survive it
\cite{ZengLiuRatnavelu2026}.  Here the certificate is probabilistic rather
than algebraic, and its scope remains the family defined next.

For \(s>0\) and \(z\in\Z^n\), write
\[
 D_{\Z^n,s}(z)=
 \frac{\exp(-\pi\lVert z\rVert_2^2/s^2)}{\rho_s(\Z^n)},
 \qquad
 \rho_s(\Z^n)=\sum_{w\in\Z^n}
 \exp(-\pi\lVert w\rVert_2^2/s^2).
\]
Let \(\nu_\gamma\) denote the centered real Gaussian distribution with
variance \(\gamma^2/(2\pi)\).  A function is \emph{negligible} if, for
every constant \(C>0\), it is eventually smaller than \(n^{-C}\).

Throughout, \(\log\) is the natural logarithm.  Set
\[
 \epsilon=2^{-n},\quad L=\Z^n,\quad A=I_n,\quad
 r=k=n^2,\quad \alpha=\frac12,\quad m=3n^5,
 \tag{1}
\]
and let \(M\) be an arbitrary integer satisfying
\[
 M>mF(n),
 \qquad
 F(n)=2^{\,n\log\log n/\log n}.
 \tag{2}
\]
Put \(c=2M\one\), \(\delta=e_1/4\), and \(y=c+\delta\).
Substitution in the sample formula of Dean--Goldsmith Lemma~7 gives the
constructed law \(P_{n,M}\):
\[
 a\sim D_{\Z^n,n^2},\qquad
 b=2\sum_{j=1}^n a_j+\frac{a_1}{4M}+e,
 \qquad e\sim\nu_{1/2},
 \tag{3}
\]
with \(a\) and \(e\) independent.

A legal comparison law \(Q_{n,M,s,\beta}\) has a fixed secret
\[
 s\in{[0,M)}^n\cap\Z^n,
 \qquad 0<\beta\le\frac12,
 \tag{4}
\]
and samples
\[
 a\sim D_{\Z^n,n^2},\qquad
 b=\frac{\langle a,s\rangle}{M}+e_\beta,
 \qquad e_\beta\sim\nu_\beta,
 \tag{5}
\]
again independently.  Both laws below are product laws across samples.

For \(N=n^8\), define
\[
 X_i=\langle a_i,\one\rangle,\qquad
 D=\sum_{i=1}^N X_i^2,\qquad
 \widehat\theta=\frac{\sum_{i=1}^N X_i b_i}{D}
 \quad(D>0).
 \tag{6}
\]
The test declares \emph{constructed} precisely when
\(D>0\) and \(\widehat\theta>3/2\); it declares \emph{legal}
otherwise.

\begin{theorem}[Explicit-family separator]\label{thm:main}
For every sufficiently large integer \(n\), every integer \(M\)
satisfying \emph{(2)}, and the data in \emph{(1)}, the target
\(y=2M\one+e_1/4\) satisfies the smoothing, parameter, constellation,
effective-noise, BDD-distance, and uniqueness promises used by the
Dean--Goldsmith Lemma~7 sample construction.

Moreover, the test in \emph{(6)} uses \(O(Nn)\) arithmetic operations
and, for every such \(M\), obeys
\begin{align*}
 \Prob_{P_{n,M}^{\otimes N}}(\text{test declares legal})
   &\le \mu(n),\\
 \sup_{\substack{s\in{[0,M)}^n\cap\Z^n\\0<\beta\le1/2}}
 \Prob_{Q_{n,M,s,\beta}^{\otimes N}}
   (\text{test declares constructed})
   &\le \mu(n),
\end{align*}
where
\begin{align}
 \mu(n)={}&2n^9\exp(-\pi n^2)
 +2n^2\exp\!\left(
 -\frac{\exp(-2\pi)}{1{,}280{,}000}n^2\right)
 \notag\\
 &\quad+2\exp\!\left(
 -\frac{49\pi\exp(-\pi)}{1600}n^{13}\right).
 \tag{7}
\end{align}
In particular, \(\mu\) is negligible, so the separation holds with
overwhelming probability uniformly over all legal laws in \emph{(4)}.
\end{theorem}

\section{Audit of the construction promises}

\begin{proposition}\label{prop:promises}
The deterministic assertion in the first paragraph of
Theorem~\ref{thm:main} holds.
\end{proposition}

\begin{proof}
For \(L=\Z^n\), the dual lattice is \(\Z^n\) and
\(\lambda_1(L)=1\).  The least eigenvalue parameter of \(A^TA\) is
\(\sigma=1\).  Claim~2 of \cite{DeanGoldsmith2017} gives
\[
 \eta_{2^{-n}}(\Z^n)\le\sqrt n.
\]
Consequently \(r=n^2>\sqrt2\,\eta_{2^{-n}}(L)\) for \(n\ge2\).
The other displayed parameter inequality used by the construction is
\[
 \frac{m\alpha}{k^2}
 =\frac{3n^5/2}{n^4}
 =\frac32n>\sqrt n.
 \tag{8}
\]
All of \(r,k,m\) are positive, \(0<\alpha<1\), and (2) is exactly
the required constellation lower bound; it also implies \(M\ge2\)
eventually.  The scaled independent response noise has parameter
\(k\alpha/r=1/2\le\alpha\).  On the orthogonal lattice used here,
the discrete Gaussian is a product of one-dimensional discrete
Gaussians, so the sample required by the constructor is available in
its standard large-width regime.

The BDD offset has norm \(\lVert\delta\rVert_2=1/4<r\).  Choose the
fixed decoding radius \(d_0=3/8\).  The distance ceiling in Lemma~7 is
\[
 \frac{M\sigma\alpha}{k^2r\sqrt2}
 =\frac{M}{2\sqrt2\,n^6}
 >\frac{3}{2\sqrt2}\frac{F(n)}{n}.
 \tag{9}
\]
Now
\[
 \log F(n)=(\log2)\frac{n\log\log n}{\log n},
\]
which dominates \(C\log n\) for every fixed \(C\).  Hence
\(F(n)/n\to\infty\), and the right side of (9) eventually exceeds
\(d_0\).

Finally, \(c=2M\one\in\Z^n\) and
\[
 \dist(y,\Z^n)=\lVert y-c\rVert_2=\frac14<d_0<\frac12
 =\frac{\lambda_1(L)}2.
\]
Coordinatewise nearest-integer inspection shows that \(c\) is the unique
closest lattice vector.  Thus the target meets the distance and
uniqueness promises.  Substituting \(A=I_n\) and \(k=r=n^2\) into
the constructor gives (3), so the audit concerns the stated law rather
than a modified one.
\end{proof}

\section{Discrete-Gaussian estimates}

Let \(r=n^2\) and let \(Z\sim D_{\Z,r}\).  Its coordinates are
centered, symmetric, and independent in \(D_{\Z^n,r}\).  Write
\[
 V=\E Z^2,\qquad v_0=\frac{\exp(-\pi)}8.
\]

\begin{lemma}[Moment and tail bounds]\label{lem:dg}
For every real \(t\) and every \(u\ge0\),
\begin{align}
 \E\exp(tZ)&\le
   \exp\!\left(\frac{r^2t^2}{4\pi}\right),
 \tag{10}\\
 \Prob(|Z|>u)&\le
   2\exp\!\left(-\frac{\pi u^2}{r^2}\right).
 \tag{11}
\end{align}
For integer \(r\ge1\),
\[
 V\ge v_0r^2.
 \tag{12}
\]
\end{lemma}

\begin{proof}
Completing the square gives
\[
 \E\exp(tZ)
 =\exp\!\left(\frac{r^2t^2}{4\pi}\right)
 \frac{\rho_r(\Z-r^2t/(2\pi))}{\rho_r(\Z)}.
\]
Poisson summation writes a shifted Gaussian mass as
\[
 \rho_r(\Z-x)
 =r\sum_{k\in\Z}\exp(-\pi r^2k^2)\exp(2\pi i kx).
\]
The triangle inequality bounds this nonnegative real quantity by its
value at \(x=0\).  This proves (10).  Chernoff's bound on each tail,
optimized at \(t=2\pi u/r^2\), proves (11).

For (12), monotonicity and integration yield
\[
 \rho_r(\Z)=1+2\sum_{k\ge1}\exp(-\pi k^2/r^2)
 \le1+r\le2r.
\]
There are at least \(r/2\) positive integers in \([r/2,r]\).
Counting these integers and their negatives shows that the unnormalized
second moment is at least
\[
 2\left(\frac r2\right)\left(\frac{r^2}{4}\right)\exp(-\pi)
 =\frac{\exp(-\pi)}4r^3.
\]
Division by \(\rho_r(\Z)\le2r\) gives (12).
\end{proof}

For independent public vectors \(a_1,\ldots,a_N\), define their
empirical second-moment matrix
\[
 C=\frac1N\sum_{i=1}^N a_i a_i^T.
 \tag{13}
\]
Thus \(\E C=VI_n\) and \(D=N\one^TC\one\).

\begin{lemma}[A common covariance event]\label{lem:cov}
For \(N=n^8\) and \(n\ge3\),
\[
 \lVert C-VI_n\rVert_{\mathrm{op}}\le\frac V{50}
 \tag{14}
\]
except on an event of probability at most
\[
 \mu_C(n)=2n^9\exp(-\pi n^2)
 +2n^2\exp\!\left(
 -\frac{\exp(-2\pi)}{1{,}280{,}000}n^2\right).
 \tag{15}
\]
\end{lemma}

\begin{proof}
Set \(T=rn\).  For \(1\le p,q\le n\), truncate the product in
entry \((p,q)\) by
\[
 W_i^{pq}=a_{ip}a_{iq}
 \mathbf1_{\{|a_{ip}|\le T,\ |a_{iq}|\le T\}}.
\]
Its range has length at most \(2T^2\).  Tail integration applied to
(11) gives
\begin{align*}
 \E[Z^2;|Z|>T]
 &\le 2T^2\exp(-\pi n^2)
   +4\int_T^\infty x\exp(-\pi x^2/r^2)\,dx\\
 &\le3r^2n^2\exp(-\pi n^2),\\
 \E[|Z|;|Z|>T]
 &\le r\left(2n+\frac1{\pi n}\right)\exp(-\pi n^2),
 \qquad \E|Z|\le r.
\end{align*}
Independence and symmetry therefore imply, for every \(p,q\),
\[
 \left|\E W_i^{pq}-V\mathbf1_{p=q}\right|
 \le6r^2n^2\exp(-\pi n^2).
 \tag{16}
\]
For \(n\ge3\), (12) makes the right side at most \(V/(100n)\).
Indeed, it is enough that
\(4800n^3\exp[-\pi(n^2-1)]\le1\), which holds at \(n=3\) and
decreases thereafter.

Put \(u=V/(100n)\).  Hoeffding's inequality
\cite{Hoeffding1963} and the range bound give
\begin{align}
 \Prob\left(\left|\frac1N\sum_iW_i^{pq}
 -\E W_i^{pq}\right|>u\right)
 &\le2\exp\!\left(-\frac{Nu^2}{2T^4}\right)\notag\\
 &\le2\exp\!\left(
 -\frac{N\exp(-2\pi)}{1{,}280{,}000n^6}\right).
 \tag{17}
\end{align}
By (11), the probability that any of the \(Nn\) sampled coordinates
has absolute value greater than \(T\) is at most
\(2Nn\exp(-\pi n^2)\).  Outside that event, truncated and original
empirical products agree.  Union-bound (17) over the \(n^2\) pairs
and use (16).  With \(N=n^8\), every entry of \(C-VI_n\) then has
absolute value at most \(V/(50n)\), with failure probability (15).
The matrix is symmetric, so its operator norm is at most its maximum
absolute row sum, which is \(V/50\).
\end{proof}

\section{Uniform separation}

We now prove the probabilistic part of Theorem~\ref{thm:main}.  Work
first on the event (14), and write \(C=VI_n+E\).  The denominator in
(6) satisfies
\[
 D=N\one^TC\one\ge\frac{49}{50}NVn>0.
 \tag{18}
\]
Thus every division below occurs only after a deterministic lower bound
on the empirical denominator.

For a legal secret, put \(t=s/M\in{[0,1)}^n\).  Its noiseless
empirical slope is
\[
 R_Q=\frac{\one^TCt}{\one^TC\one}.
\]
Since \(\one^Tt\le n\), \(\lVert t\rVert_2\le\sqrt n\), and
\(\lVert E\rVert_{\mathrm{op}}\le V/50\),
\[
 R_Q\le
 \frac{Vn+(V/50)\sqrt n\sqrt n}{Vn-(V/50)n}
 =\frac{51}{49},
 \qquad
 \frac32-R_Q\ge\frac{45}{98}>\frac14.
 \tag{19}
\]
This is simultaneous for every \(t\) in the full legal cube.

Under the constructed law, the noiseless empirical slope is
\[
 R_P=2+\frac1{4M}
 \frac{\one^TCe_1}{\one^TC\one}.
\]
The same operator-norm event gives
\[
 \left|\frac{\one^TCe_1}{\one^TC\one}\right|
 \le\frac{V+(V/50)\sqrt n}{(49/50)Vn}
 =\frac{50+\sqrt n}{49n}\le1
 \tag{20}
\]
for \(n\ge2\).  Since \(M\ge1\),
\[
 R_P\ge\frac74,
 \qquad R_P-\frac32\ge\frac14.
 \tag{21}
\]

It remains to control the response noise.  Conditional on the public
vectors, if the independent errors have law \(\nu_\gamma\), then for
\(D>0\)
\[
 \frac{\sum_iX_ie_i}{D}
 \sim N\!\left(0,\frac{\gamma^2}{2\pi D}\right).
 \tag{22}
\]
Here \(\gamma=1/2\) for \(P_{n,M}\) and
\(0<\gamma=\beta\le1/2\) for every legal law.  A two-sided Gaussian
tail bound, (19), and (21) show that either classification error,
conditional on (14), is at most
\[
 2\exp(-\pi D/4)
 \le2\exp\!\left(
 -\frac{49\pi\exp(-\pi)}{1600}n^{13}\right),
 \tag{23}
\]
where (12), (18), \(r=n^2\), and \(N=n^8\) were used.  Adding
(15) proves (7).  The constants do not depend on \(M\), \(s\), or
\(\beta\).  Every term in (7) is a polynomial factor times
\(\exp(-cn^2)\), or \(\exp(-cn^{13})\), for a fixed \(c>0\);
hence their sum is negligible.

If \(D=0\), the prescribed convention declares the sample legal.  Under
the constructed law, (18) shows that this event is contained in the
covariance failure event.  Under a legal law it cannot cause a false
declaration of constructed.  This covers the zero denominator and
completes the proof of Theorem~\ref{thm:main}.

\section{What the result does and does not show}

The population version of the covariance ratio already exposes the
mechanism.  With \(X=\sum_j a_j\) and \(V=\E a_1^2\),
\[
 \frac{\E_{P}[Xb]}{\E[X^2]}
 =2+\frac1{4Mn}>2,
 \qquad
 \frac{\E_{Q}[Xb]}{\E[X^2]}
 =\frac1{Mn}\sum_js_j<1.
\]
The proof above is needed because a ratio of expectations does not
control the random empirical ratio.  Lemma~\ref{lem:cov} supplies one
event that uniformly controls every legal secret, keeps the denominator
positive, and retains a constant margin under both hypotheses;
conditioning then handles the Gaussian response noise.

This advances Li's exact distributional mismatch for the single family
(1)--(3) to an explicit polynomial-sample, overwhelming-probability
separator.  It does not show that every invalid oracle call in the
Dean--Goldsmith reduction is efficiently detectable, does not give a
decoder for arbitrary MIMO instances, and does not establish or refute
the security of physical-layer cryptography in general.

\section*{Computational and AI assistance disclosure}

OpenAI Codex assisted with literature navigation, adversarial checking of
the argument, algebraic organization, and manuscript preparation.  A
separate model context reconstructed and audited the stated mathematical
claim before manuscript preparation; that audit is not formal
verification or peer review.  No numerical experiment, computer-algebra
output, or floating-point calculation is used as proof: the complete
certificate is the analytic argument and explicit inequalities above.
The author is responsible for every claim and for compliance with any
venue's disclosure and authorship policies.

\bibliographystyle{alpha}
\bibliography{references}

\end{document}
